% ==============================================================================
% ================================ Zusammenfassung ==============================
% ==============================================================================

%\chapter{Zusammenfassung und Ausblick}
%\label{chap:zusammenfassung}
%\thispagestyle{fancy}

\section{Zusammenfassung und Ausblick}

Die vorliegende Arbeit hat gezeigt, dass die effektive Visualisierung von Sensordaten aus Wearables einen entscheidenden Beitrag zur Früherkennung chronischer Herzerkrankungen leisten kann \textit{(\acs{vgl.} \cite[4]{2})}. Ein zentrales Ergebnis ist, dass die Transformation von Rohdaten in zielgruppengerechte Informationen nur durch einen strukturierten Prozess der Datenvorbereitung – von der Bereinigung bis zum Mapping – gelingen kann \textit{(\acs{vgl.} \cite[11, 31]{1})}. Dabei hat sich gezeigt, dass unterschiedliche Nutzergruppen spezifische Anforderungen stellen: Während Medizinerinnen und Mediziner zur Mustererkennung auf hochdimensionale Darstellungen wie die Streudiagramm-Matrix angewiesen sind, benötigen Patientinnen und Patienten intuitive Formate wie Liniendiagramme zur Förderung der Therapietreue \textit{(\acs{vgl.} \cite[7]{2})}.

Ein weiterer wesentlicher Aspekt ist die ethische Verantwortung bei der Gestaltung von Datenvisualisierungen. Manipulationen durch abgeschnittene Achsen oder verzerrende 3D-Effekte müssen konsequent vermieden werden, um Fehlinterpretationen im medizinischen Kontext auszuschließen \textit{(\acs{vgl.} \cite[50]{1}, \cite[43]{2})}. Die praktische Umsetzung in R hat zudem verdeutlicht, dass moderne Werkzeuge wie RMarkdown eine effiziente und reproduzierbare Kombination aus Analyse, Visualisierung und Storytelling ermöglichen, um auch Laien komplexe Zusammenhänge verständlich zu vermitteln \textit{(\acs{vgl.} \cite[1, 55]{3})}.

\subsection{Ausblick}

Zukünftige Erweiterungen des Visualisierungskonzepts könnten die Integration von Echtzeit-Dashboards mithilfe von Paketen wie \texttt{shiny} umfassen, um eine kontinuierliche Live-Überwachung der Patientendaten zu ermöglichen \textit{(\acs{vgl.} \cite[49]{3})}. Darüber hinaus bietet die Analyse höherdimensionaler Datenräume \textit{(\(E_n^P\) mit \(n > 6\))} Potenzial für noch präzisere Diagnosemodelle \textit{(\acs{vgl.} \cite[33]{1})}. Angesichts der stetig wachsenden Datenmengen wird die Bedeutung einer klaren und ehrlichen visuellen Kommunikation weiter zunehmen, um der Informationsüberflutung entgegenzuwirken und fundierte medizinische Entscheidungen zu unterstützen \textit{(\acs{vgl.} \cite[43, 52]{2})}.

% ==============================================================================
% ================================ Zusammenfassung ==============================
% ==============================================================================